\section{Arithmetic comparison of the histories}\label{sec:arithmetic}

We prove Proposition~\ref{prop:comparisons}. Throughout this section all the
constants in the construction and the depth $k$ are fixed, and then $L$ tends
to infinity. The notation, priors, templates, and support conventions are
those of Section~\ref{sec:construction}. At comparison level $l$ put $r=2^l$.
We treat both the final environment and the diagonal environment of
Proposition~\ref{prop:comparisons}; in the latter the positive giant is the
external integer variable. All estimates below are uniform in the assignments
being compared and in the permitted frequency histories.

The arithmetic reduction first integrates the top giants at the regular
and internal small primes. After resolving the resulting internal-prime
conditions, we replace the bulk primes by real log coordinates and independent
unit residues, keeping the real weights and frequency tests. The paired
estimate \eqref{eq:src56} then uses signed spectator correlations. For the
norm estimate \eqref{eq:src55}, we may take absolute values termwise, but
must retain frequency divisibility to control the sum over histories.

\subsection{Histories, supports, and real-variable weights}

Expand the two amplitudes in a comparison into their frequency histories.
For the norm estimate \eqref{eq:src55}, use the same assignment in the two
histories. At level $t\geq1$ there are $2^{l-t}$ nodes in a history, and its
sampled compensation vector has $4\cdot2^{t-1}$ entries. Thus the total
number of internal prime samples in the two histories is $O(kr)$. This counts
each sample when it is made, rather than each subsequent occurrence of that
sample in a descendant list.

Partition these samples according to their equality pattern. There are
$O_k(1)$ patterns. Entries in a single compensation vector must be distinct,
and the disjoint size ranges imply that equality is possible only between
samples of the same compensation type. For each distinct internal sampled
prime $b$, let $n_b$ be its multiplicity in the pattern. The product of the
weights in \eqref{eq:src52} is then
\[
                         \prod_b b^{n_b}.
\]
We retain all the other support restrictions for the moment.

Write $x_+,x_-$ for the top giant entries. The negative giant $x_-$ is an
actual prime in both environments; $x_+$ is a prime in the final environment
and an external positive integer in the diagonal environment. Let $d$ be the
product of the $m$ distinct spectator primes, and let $c_{\rm sm}$ be the
product of the distinct actual regular small primes at the top. The internal
samples are disjoint from these primes by their types. Define
\begin{equation}\label{eq:arith-frequency-modulus}
 R_f=\prod_{\nu}|\nu|,\qquad Q_f=R_f^{k+2},
 \qquad \log Q_f=O_k(m),
\end{equation}
where the product includes every frequency in the two histories. Repetitions
in this product are allowed. A modulus equal to one has its usual trivial
interpretation. Every small prime in the construction is larger than every
frequency and hence is coprime to $Q_f$.

At a reversing node write the two current giant entries as $X_+,X_-$ and
factor its child products as
\begin{equation}\label{eq:arith-child-products}
 H_+=X_+C_+M_+,\qquad H_-=X_-C_-M_-.
\end{equation}
Here $M_\pm$ are the products of the bulk slots in the respective child
subtrees, and $C_\pm$ contain their other regular small factors. The latter
may include compensation samples from preceding reversals on the path.
Both current giants are rational linear functions of $x_+,x_-$, by
\eqref{eq:src51}.

We first justify an exhaustive residue description of the support. Assume
temporarily that the top giants are coprime and that both are units modulo
$c_{\rm sm}d$. Suppose that a current state is pairwise coprime. At its
reversal the numerator is
\begin{equation}\label{eq:arith-numerator}
             N=vH_--wH_+,\qquad p=\frac{N}{su}.
\end{equation}
Every prime divisor of $H_+$ is coprime to $H_-$ and to $v$: for an actual
small prime this follows by size, and for a giant factor it is one of the
giant--frequency unit conditions. Hence $N$ is a unit at every prime factor
of $H_+$. The same argument, interchanging the two sides, applies to $H_-$.
If $p$ is integral, it follows from $N=sup$ that both $p$ and $u$ are
automatically coprime to every inherited slot in $H_+H_-$. In particular, a
new compensation sample cannot divide an inherited giant. To obtain a
pairwise-coprime child state, it remains only to require that $p$ be a unit
against its own $u$ and against $d$.

The following tests therefore describe all the remaining arithmetic support:
\begin{enumerate}
\item Modulo frequency factors, impose $s\mid N$ at every node and impose
all giant--frequency unit conditions. These tests are determined by the top
and small coordinates modulo $Q_f$. Division by a node frequency loses at
most one factor $R_f$ of precision; the $u$'s are units modulo $Q_f$. At most
$k$ successive divisions occur, so \eqref{eq:arith-frequency-modulus} leaves
enough precision for every subsequent test.
\item For each $b\mid u$ at a node, require
\begin{equation}\label{eq:arith-internal-tests}
                        N\equiv0\pmod b,
             \qquad N\not\equiv0\pmod{b^2}.
\end{equation}
The entries of the current $u$ are distinct and are coprime to $s$, so these
are exactly its remaining integrality and own-$u$ unit requirements.
Ancestor denominators involve other compensation types and frequencies;
they are units even modulo $b^2$. Thus the tests can also be computed directly
from the rational expressions in the top giants.
\item Require the top giant units modulo $c_{\rm sm}d$. At each spectator
$q\mid d$, use the zero convention in \eqref{eq:src36} whenever an inserted
giant vanishes modulo $q$.
\end{enumerate}
These tests, together with positivity and the real cutoffs, agree with
ordinary recursive evaluation over the integers. Indeed, assuming the
ancestors have already been evaluated integrally, the first two tests give
divisibility by the relatively prime factors $s$ and $u$. The preceding
coprimality argument then passes the required support to the children. This
proves the assertion by induction down the tree.

At a spectator $q$, put $D'=d/q$. Formula
\eqref{eq:arith-child-products}, the order of the child giants, and
\eqref{eq:src51} identify the spectator factors from a history with the tree
value in \eqref{eq:src36}. In particular, if the bulk product at leaf $i$ is
$M_i$, the constants at a child have the product consistency required in that
formula. Denote the two tree values by $W_{1,q},W_{2,q}$. Their combined
spectator factor is
\begin{equation}\label{eq:arith-spectator-product}
                         \prod_{q\mid d}W_{1,q}\overline{W_{2,q}}.
\end{equation}
The zero convention incorporates precisely the inserted-giant unit tests at
$q$.

Let $\Xi$ denote the product of all smooth real factors in the two histories,
including $K_{\rm ar}$ when it is present, but excluding
$\prod_b b^{n_b}$. Define it for real positive top and small variables by
using \eqref{eq:src51} over $\R$. Set the contribution to zero if an inserted
$p$ is nonpositive. This is a smooth extension: the factor
$\varphi(\log p-G)$ already vanishes for $p<e^{G-1}$.
Counterpart role indicators involving nonbulk slots are retained as
restrictions on those variables; all bulk priors and their role ranges are
identical, so no such indicator cuts a bulk log coordinate.

There are $r$ leaf factors in each history. On their common support, the
modulus in each instance of \eqref{eq:src50} has logarithm
$\log X+\Delta_0+O_k(1)$. This follows from the giant cutoffs, all the
specified small-prime ranges, and both sum bins at that leaf. Consequently
\begin{equation}\label{eq:arith-xi-bound}
                         |\Xi|\leq e^{-r\Delta_0+O_k(1)}.
\end{equation}
Its first derivatives in the log coordinates of the top giants and bulk
slots are bounded by $\exp(O_k(m))$. To check the possible cancellation in
\eqref{eq:arith-numerator}, use \eqref{eq:src48} at each subtree. Each of
$|vH_-|$ and $|wH_+|$, divided by $|s|up$, is at most
\[
                         \exp(\Delta_j+E_{j-1}+O_k(1))
\]
at a node of level $j$. The depth is bounded by $k$, so repeated logarithmic
differentiation along a path costs $\exp(O_k(m))$. The remaining cutoffs,
the bounded Fourier arguments in the leaf factors, and $K_{\rm ar}$ satisfy
the same bound. The bounds extend over support boundaries: any problematic
nonpositive inserted value lies outside the giant cutoff, where the smooth
extension is zero. We use the actual prior ranges as integration domains,
so no extra discontinuous real indicator is needed.

\subsection{A progression estimate retaining the exceptional term}

We record the precise approximation needed for the two successive
idealizations. In the following table, either row may be used:
\begin{equation}\label{eq:arith-progression-scales}
\begin{array}{c|rrrrrr}
 &a_0&a_1&\mu&\delta&\delta'&\theta\\ \hline
 \mathrm{giant}&.049&.95&.012&.015&.018&.022\\
 \mathrm{bulk}&.0039&.007&.0012&.0014&.0016&.0018
\end{array}
\end{equation}

\begin{lemma}\label{lem:arith-progression}
Fix a row of \eqref{eq:arith-progression-scales}. There is at most one
primitive real character $\chi_*$, with conductor at most
$\exp(\exp(\mu L))$, and one associated real zero $\beta_*<1$ which must
be retained in the following formula. If
\[
 \log M_*\leq e^{\mu L},\qquad (a,M_*)=1,
\]
and $I$ is an interval of length at most one in
$[e^{a_0L},e^{a_1L}]$, then
\begin{equation}\label{eq:src57}
 \sum_{\substack{p\ \text{prime}\,,\ \log p\in I\\p\equiv a\pmod{M_*}}}
       \frac1p
 =\frac1{\phi(M_*)}\int_I
       \bigl(1-\chi_*(a)e^{(\beta_*-1)t}\bigr)\frac{dt}{t}
   +O\bigl(\exp(-c e^{\delta'L})\bigr).
\end{equation}
The exceptional term is omitted if no such character exists or its conductor
does not divide $M_*$. The implied constant is absolute for the fixed row.
The multiplier inside the integral lies in $[0,2]$.
\end{lemma}

\begin{proof}
Put $Q=\exp(e^{\mu L})$ and $S=\exp(e^{\theta L})$. The classical
zero-free region and Page's theorem give, simultaneously for primitive
conductors up to $Q$ and ordinates of absolute value at most $S$,
\begin{equation}\label{eq:arith-zero-free}
                     1-\Re\rho\geq c_0e^{-\theta L},
\end{equation}
apart from at most one simple real zero of a primitive real character.
One may take for $\chi_*$ the character of a zero in this latter narrow
region, if there is one. This follows from
\cite[Theorem~11.3 and Corollary~11.10]{MontgomeryVaughan2007}; the analogous
zero-free statement for $\zeta$ has no exception. We also use the standard
zero count $N(T,\chi)\ll T\log(q(T+2))$ for primitive conductor $q$;
see \cite[Section~3.7.3]{Helfgott}.

Here are details giving the required short-interval precision. Let $t_0$ be
the left endpoint of $I$, put $y=e^{t_0}$, and set
$\epsilon=\exp(-e^{\delta'L})$. Sandwich the indicator of the corresponding
interval in the variable $u=n/y$ between smooth nonnegative functions
$f_-,f_+$ supported in $[1/2,4]$. They may be chosen with
\[
 \int(f_+-f_-)\,du=O(\epsilon),\qquad
 \|f_\pm^{(j)}\|_\infty\ll_j\epsilon^{-j}.
\]
This remains possible when the interval is shorter than $\epsilon$, by
taking the lower function zero. For either smooth function let
\[
 F(u)=\frac{f(u)}{u\log(yu)},\qquad
 \mathcal M F(s)=\int_0^\infty F(u)u^{s-1}\,du.
\]
For a primitive character $\chi$, Mellin inversion and a contour shift give
\begin{equation}\label{eq:arith-smoothed-explicit}
 \frac1y\sum_{n\geq1}\Lambda(n)\chi(n)F(n/y)
 =\ind_{\chi=1}\int_0^\infty F(u)\,du
   -\sum_\rho y^{\rho-1}\mathcal M F(\rho)
   +O\bigl(y^{-1+o(1)}\bigr).
\end{equation}
The sum is over nontrivial zeros. For completeness, shift the Mellin integral
of $-L'/L$ to real part $-1/2$, taking limits through heights avoiding zeros.
The functional equation bounds the logarithmic derivative on that line by
$O(\log(q(|t|+2)))$. Mellin decay makes the shifted integral convergent and
bounded by a fixed power of $\epsilon^{-1}$ times
$y^{-3/2}\log(q+2)$; a possible residue at zero costs
$y^{-1}$ times such a factor. Because $a_0>\delta'$, these bounds have the
stated size. The pole and zero residues are exactly those displayed in
\eqref{eq:arith-smoothed-explicit}. This is the usual explicit-formula
argument underlying \cite[Theorem~11.16]{MontgomeryVaughan2007}.

On $0\leq\Re s\leq1$, integration by parts gives, for fixed $j$,
\begin{equation}\label{eq:arith-mellin-decay}
 |\mathcal M F(s)|\ll1,\qquad
 |\mathcal M F(s)|\ll_j
       \epsilon^{-O(j)}(1+|\Im s|)^{-j}.
\end{equation}
The nonexceptional zeros of height at most $S$ therefore contribute at most
\begin{equation}\label{eq:arith-low-zero-error}
 \exp\bigl(-c e^{(a_0-\theta)L}+e^{\theta L}+O(L)\bigr).
\end{equation}
For the giant row $a_0-\theta=.027>.022=\theta>\delta'$, and for the bulk
row $a_0-\theta=.0021>.0018=\theta>\delta'$. Thus
\eqref{eq:arith-low-zero-error} is smaller than the required error. For
zeros above $S$, dyadic summation of the zero count and
\eqref{eq:arith-mellin-decay} gives
\[
 \exp\bigl(C_j e^{\delta'L}-(j-1)e^{\theta L}+O(L)\bigr),
\]
which is also sufficient, since $\delta'<\theta$.

For a prime $n=p$, the summand $\Lambda(p)F(p/y)/y$ is $f(p/y)/p$.
Prime powers of higher exponent give $y^{-1/2+o(1)}$. Passing from an
induced character modulo $M_*$ to its primitive character changes only the
prime powers at prime divisors of $M_*$ and gives a negligible error as
well. Orthogonality of the characters modulo $M_*$ now selects the residue
class $a$; its factor $1/\phi(M_*)$ cancels the number of character errors.
The possible exceptional primitive character occurs among these induced
characters precisely when its conductor divides $M_*$.

Under $u=e^{t-t_0}$, the principal integral becomes $f(e^{t-t_0})\,dt/t$.
The exceptional zero term becomes
\[
              f(e^{t-t_0})e^{(\beta_*-1)t}\frac{dt}{t}.
\]
In particular no factor $\beta_*$ is left over. Equivalently, this follows
by differentiating the exceptional Chebyshev term
$-\chi_*(a)y^{\beta_*}/(\beta_*\phi(M_*))$; see
\cite[Corollaries~11.17 and~11.20]{MontgomeryVaughan2007}.
The resulting density is nonnegative and at most $2\,dt/(\phi(M_*)t)$.
The upper and lower smooth tests therefore differ in their main integrals
by $O(\epsilon)$, and their pointwise sandwich on the positive prime measure
proves \eqref{eq:src57}.
\end{proof}

For each row, choose the possible exceptional character in
Lemma~\ref{lem:arith-progression} before fixing the nongiant priors. If its
conductor has a prime factor at least
\begin{equation}\label{eq:arith-deletion-threshold}
                         T_*=\exp(e^{.0004L}),
\end{equation}
delete one such prime value from every nongiant prior. This prescribes at
most two deletions, as allowed in the construction. Every prime factor of
the moduli below outside $Q_f$ exceeds $T_*$, whereas every prime factor of
$Q_f$ is smaller than $T_*$ by \eqref{eq:arith-frequency-modulus}. Hence if
an exceptional conductor still divides one of these moduli, its full
conductor divides $Q_f$. Indeed, a selected large prime factor has been
excluded, and otherwise all conductor prime factors must lie in the $Q_f$
part; coprimality with the other parts also accounts for their exponents.
Thus every surviving Page correction depends only on the real coordinate
and the $Q_f$ residue. It does not couple any of the other CRT coordinates.

\subsection{Idealization of the top giants}

Condition on all actual small primes, including the internal samples in
both histories. At fixed frequencies, the logarithm of the product of their
large weights and all pointwise transform bounds is at most $e^{.012L}$ for
large $L$. More explicitly, the regular and internal small-prime products
have logarithms
\begin{equation}\label{eq:arith-small-size-budget}
 O_k(e^{.01L})+O_k(m e^{.006L}),
\end{equation}
and the spectator products have still smaller logarithms. Use
$|g_\ell|\leq\sqrt\ell$ and $|\widetilde g_\ell|\leq1$ for their local
pointwise bounds. All the factors $\prod b^{n_b}$ are included in this
estimate.

We may now remove the temporary assumption that the two top giants are
coprime, using the residue description just obtained as the extension.
In the prime environment the discrepancy is equality of the two primes;
its probability is at most $e^{-G+O(L)}$. In the external-integer
environment, a fixed prime $x_-$ has only $O(1)$ positive multiples in the
log cell for $x_+$, and their normalized total mass is $O(e^{-G})$.
Multiplying by the preceding pointwise bound is harmless, because
$G\geq e^{.049L}$.

Use the modulus
\begin{equation}\label{eq:arith-giant-modulus}
                         M_*=c_{\rm sm}dQ_f\prod_b b^2.
\end{equation}
Its factors are pairwise coprime, and
\eqref{eq:arith-small-size-budget} gives $\log M_*\leq e^{.012L}$.
The giant log cells lie in $[e^{.049L},e^{.95L}]$, by their selection from
\eqref{eq:src28}. Lemma~\ref{lem:arith-progression}, with the original
cell cutoff and normalization, replaces each top prime by a real log
coordinate and a Haar-unit residue modulo $M_*$. Its density is the
principal harmonic density times the correction in \eqref{eq:src57}.
For an external integer the analogous replacement uses
\begin{equation}\label{eq:arith-external-density}
                         e^{t-G}\varphi(t-G)\,dt
\end{equation}
and uniform residues on all classes modulo $M_*$. Elementary counting of
integers in a progression gives this replacement with error
$e^{-G+O(L)}$ per interval and residue. The total ideal measure of each
coordinate is bounded. A prime cell's normalizing reciprocal costs only
$\exp(O(L))$.

We give a joint error estimate so that no conditional equidistribution is
implicit. Partition both log cells using mesh
$\eta_g=\exp(-e^{.015L})$ and freeze the smooth factor $\Xi$ in each box.
All other tests at fixed small variables and frequencies are residue tests.
The variation error, including the pointwise costs above, is bounded by
\[
 \exp\bigl(-e^{.015L}+O(e^{.012L})+O_k(m)\bigr).
\]
The logarithm of the total number of boxes and pairs of residue classes is
$O(e^{.015L}+e^{.012L})$. Thus the absolute interval errors from the giant
row of \eqref{eq:src57}, even summed over all these boxes and classes and
multiplied by the pointwise bound, contribute at most
\[
 \exp\bigl(-c e^{.018L}+O(e^{.015L})+O(e^{.012L})+O_k(m)\bigr).
\]
The integer-coordinate errors have the still stronger negative term $-G$.
It follows that the whole giant replacement has uniform error
\begin{equation}\label{eq:arith-giant-error}
                         O\bigl(\exp(-e^{.013L})\bigr).
\end{equation}
The Page factors are part of the measures in this argument; they need not
be frozen or differentiated.

\subsection{Integration of the actual and internal small-prime coordinates}

In the ideal giant distribution, CRT separates its coordinates, and a
surviving Page factor uses only $Q_f$. First integrate the giant coordinates
modulo $c_{\rm sm}$. In the diagonal environment, $x_-$ is already a unit
there, and $x_+$ must be restricted from all classes to units. At
$\ell\mid H_{\rm sm}$ the argument of $g_\ell$ in \eqref{eq:src54} is a
fixed unit divided by $x_+x_-$. It is therefore uniform on the unit group
once both giants are units. Since $g_\ell(0)=0$ and
$\E_{a\bmod\ell}|g_\ell(a)|^2=1$, the average of its squared modulus on
units is $\ell/(\ell-1)$. Independence over $\ell$ gives
\begin{equation}\label{eq:src58}
 \frac{\phi(c_{\rm sm})}{c_{\rm sm}}
    \prod_{\ell\mid H_{\rm sm}}\frac{\ell}{\ell-1}
       =\prod_{\ell\mid u_{\rm out}}\left(1-\frac1\ell\right),
\end{equation}
where $u_{\rm out}$ is the actual outer compensation product of type $j$ in
the diagonal environment. This factor is bounded by one and is independent
of the bulk values. In the final environment the analogous integral equals
one. The support descent proved above shows that no remaining test uses
these giant coordinates modulo the actual regular small primes.

Fix now a distinct internal sampled prime $b$. Each of its occurrences
requires a numerator in \eqref{eq:arith-numerator} to vanish modulo $b$.
This numerator is a nonzero linear form in $(x_+,x_-)$ modulo $b$.
To see nonvanishing, along its ancestor path each substitution retains one
giant entry and replaces the other by a linear combination with two unit
coefficients. Its matrix is invertible modulo $b$, and the current numerator
is itself a nonzero row applied to this pair. All coefficients are rational
expressions in other small regular or internal primes and in the fixed
frequencies. No variable of the same compensation type as $b$ is needed:
that type has not yet appeared at any ancestor, and the current numerator
does not contain the current $u$. This also holds across the two histories,
and the ancestor transformations are invertible modulo $b^2$.

If the forms for all occurrences of $b$ have rank two, they have no solution
under the sampled giant residues, because $x_-$ is a unit. In rank one, their
common line has probability
\begin{equation}\label{eq:arith-line-probabilities}
 \begin{cases}
 (b-1)^{-1},&\text{both giants prime and both line coefficients nonzero},\\
 b^{-1},&\text{$x_+$ external and its coefficient nonzero},\\
 0,&\text{otherwise}.
 \end{cases}
\end{equation}
For example, in the second case each of the $b-1$ unit choices of $x_-$
determines one of the $b$ choices of $x_+$. These probabilities multiply over
distinct $b$ by CRT.

Conditional on a rank-one solution modulo $b$, a forbidden zero modulo
$b^2$ from \eqref{eq:arith-internal-tests} removes at most $1/b$ of the
uniform lifts: at least one coefficient of the form is a unit. A union bound
over its occurrences gives a relative loss $O_k(1/b)$. To justify dropping
these exclusions in the weighted expression, we first bound those weights.
For a prime in a compensation log cell, every occurrence prior $\mu_i$
satisfies $b\mu_i(b)\leq e^{CL}$, with $C$ independent of $k$ once $L$
is sufficiently large. Thus, selecting any occurrence as representative,
\[
 \left(\prod_{i=1}^{n_b}\mu_i(b)\right)
       \frac{b^{n_b}}{b-1}
 \leq 2\mu_1(b)e^{C(n_b-1)L}.
\]
Summing over the $O(kr)$ occurrences shows that the joint measure after
integrating these line conditions is dominated by independent
representative priors times
\begin{equation}\label{eq:src59}
                              \exp(CkrL).
\end{equation}
Restrictions that representatives be distinct can be discarded for this
upper bound. The same domination holds in the external-integer case since
$1/b\leq1/(b-1)$.

After \eqref{eq:src58}, the remaining large pointwise factor is the
spectator product. Since there are $2r$ leaf transforms at each spectator,
its logarithm is at most
\begin{equation}\label{eq:arith-spectator-cost}
                              O_k(m e^{.001L}).
\end{equation}
Every internal $b$ satisfies $\log b\asymp_k e^{.01L}$. The relative
$O_k(1/b)$ losses, multiplied by \eqref{eq:src59},
\eqref{eq:arith-spectator-cost}, and all $\exp(O_k(m))$ costs, are therefore
negligible. We drop all the nonvanishing-modulo-$b^2$ exclusions henceforth.

\subsection{Removing arithmetic coincidences in the line conditions}

For the signed comparison \eqref{eq:src56} we need to remove dependence of
the line probabilities on accidental bulk congruences. Fix the equality
pattern among internal samples. Regard the remaining small-prime
representatives and actual small primes as independent formal variables.
For each $b$, clear the denominators in the coefficient-zero and
two-by-two minor-zero tests for its forms. Those denominators are units
modulo $b$. We obtain integer polynomials in the fixed frequencies and the
formal variables, with no spectator variable and no variable of $b$'s
compensation type.

The degree is $O_k(m)$. This follows inductively because each coefficient in
a reversal is a product of a subtree's small slots times a frequency, and
there are only $k$ substitutions on any path. Clearing the products of
ancestor denominators and forming a two-by-two minor preserves this bound.
The same induction gives, whenever an evaluation is nonzero,
\begin{equation}\label{eq:arith-polynomial-size}
                               \log|P|\leq e^{.012L}
\end{equation}
for large $L$: the total degrees are $O_k(m)$, all small-prime logarithms
are $O_k(e^{.01L})$, and the logarithms of the frequency coefficients are
$O_k(m)$.

Replace each test by the question whether its polynomial is identically
zero over $\Q$. We quantify the error in this replacement under
\eqref{eq:src59}. If a polynomial of total degree $D$ in independent
variables is not identically zero and each variable has maximal atom at
most $\alpha$, then
\begin{equation}\label{eq:arith-polynomial-root-bound}
                             \Prob(P=0)\leq D\alpha.
\end{equation}
This is a maximal-atom variant of the Schwartz--Zippel argument.
Schwartz \cite[Lemma~1, p.~702]{Schwartz1980} gives the sharp total-degree
estimate for uniform finite sets; Zippel
\cite[\S 3.1, Theorem~1]{Zippel1979} gives related coordinate-degree
zero estimates in sparse interpolation.  Earlier random-evaluation
identity testing appears in DeMillo and Lipton \cite{DeMilloLipton1978}.
The proof below extends the uniform-finite-set argument to independent
nonuniform laws; that extension is supplied here, not imported from these
references.
Indeed, write it as a polynomial of degree $d$ in its last variable, with
nonzero leading coefficient of degree at most $D-d$. Induction bounds the
probability that this coefficient vanishes by $(D-d)\alpha$; otherwise
there are at most $d$ roots in the last variable. This proves
\eqref{eq:arith-polynomial-root-bound}.

All variables occurring here are at least bulk size. Their maximal atoms
are at most $\exp(-c e^{.004L})$, so an accidental exact zero has probability
at most $O_k(m)\exp(-c e^{.004L})$. If the evaluation is nonzero,
\eqref{eq:arith-polynomial-size} implies that it has at most $e^{O(L)}$
distinct prime divisors. The representative $b$ is independent of this
evaluation, because its whole compensation type is absent from the
polynomial. Its maximal atom is at most
$\exp(-c_k e^{.01L}+O(L))$, so the probability that it divides this
nonzero evaluation is again negligible. A polynomial which is an identity
vanishes modulo $b$ identically, because the cleared denominators are units.

The error bounds remain valid with the other support restrictions: for this
purpose discard those restrictions and use the representative-prior
domination. Both the original and the replaced line factors are bounded by
$1/(b-1)$, so the same domination applies to their difference. There are only
$O_k(1)$ tests, and multiplying by the spectator bound
\eqref{eq:arith-spectator-cost} and all $\exp(O_k(m))$ factors still gives
an error
\begin{equation}\label{eq:arith-symbolic-error}
                              O\bigl(\exp(-e^{.002L})\bigr).
\end{equation}

We now integrate out the giant coordinates modulo each internal $b$.
For \eqref{eq:src56} their contribution is its symbolic rank and feasibility
flag times the appropriate scalar in \eqref{eq:arith-line-probabilities}.
The flags depend on the formal pattern and frequencies, not on the sampled
bulk values. The line itself leaves no further condition, since no other
factor uses the giant coordinates at $b$ after its $b^2$ exclusions have
been removed. Thus replacing the flags and integrating the line conditions
removes all dependence on bulk residues modulo internal primes.

For \eqref{eq:src55} there is a simpler nonnegative upper bound. Take
absolute values in the two-history expansion and replace every line
probability by its upper bound $1/(b-1)$ before using
\eqref{eq:src59}. In this estimate rank and feasibility flags are discarded;
they do not remain as restrictions on the bulk variables. The symbolic
replacement is needed for the signed comparison, but not for this upper
bound.

We may also remove distinctness among the actual bulk samples in the
simplified expressions. There are $rm=O_k(m)$ such samples, and their
collision probability under independent priors is at most their number of
pairs times their maximal point mass. Equations~\eqref{eq:src59} and
\eqref{eq:arith-spectator-cost} show that the resulting error is bounded by
\eqref{eq:arith-symbolic-error}. This step extends the already simplified
formula; it does not attempt to use \eqref{eq:src58} with a nonsquarefree
$c_{\rm sm}$. Keep spectator distinctness and any restrictions involving
only nonbulk variables.

\subsection{Joint idealization of the bulk variables}

The remaining dependence on actual bulk values is smooth in their log
coordinates or occurs modulo $dQ_f$. No condition involving their residues
at internal sampled primes or at actual regular small primes remains.
Furthermore
\begin{equation}\label{eq:arith-bulk-modulus}
              \log(dQ_f)=O_k(m e^{.001L})<e^{.0012L}.
\end{equation}
Use the bulk row of Lemma~\ref{lem:arith-progression} to replace the $rm$
bulk priors by their real harmonic densities and Haar-unit residues modulo
$dQ_f$, retaining the possible Page corrections. The actual range
$e^{.004L}\leq\log p\leq e^{.006L}$ is strictly inside the row's permitted
range. Its harmonic normalization is of order $L$. If a deleted prime lies
in the bulk range, removing its single atom has negligible cost by the same
maximal-atom estimate just used. Thus the approximation may use the all-prime
formula, while keeping the given normalization constants.

Here too the approximation is joint. Let $N_b=rm$ and partition every bulk
log range using mesh $\eta_b=\exp(-e^{.0014L})$. The logarithm of the number
of joint boxes and joint residue classes is at most
\begin{equation}\label{eq:arith-bulk-box-count}
 O_k\bigl(m(e^{.0014L}+L+\log(dQ_f))\bigr).
\end{equation}
Freeze the smooth weight in each box. Equations~\eqref{eq:src59} and
\eqref{eq:arith-spectator-cost} bound its remaining total or pointwise
cost by $\exp(O_k(m e^{.001L}))$. The derivative bound gives variation
error at most
\[
 \exp\bigl(-e^{.0014L}+O_k(m e^{.001L})+O_k(m)\bigr).
\]
For the interval errors, even summing over all the boxes and residues in
\eqref{eq:arith-bulk-box-count} gives the upper bound
\[
 \exp\bigl(-c e^{.0016L}
       +O_k(m(e^{.0014L}+L+e^{.0012L}))
       +O_k(m e^{.001L})\bigr).
\]
Both are
\begin{equation}\label{eq:arith-bulk-error}
                              O\bigl(\exp(-e^{.00125L})\bigr).
\end{equation}
The bounded masses of the other coordinates add only $\exp(O_k(m))$.
In the main term their real distributions are integrated as measures; the
box count is charged only to the approximation errors. There is therefore
no box-count factor multiplying a main term.

For fixed $k$ the number of frequency histories is $\exp(O_k(m))$, because
each frequency has bound $\exp(O_k(m))$ and there are $O_k(1)$ nodes. The
number of equality patterns is $O_k(1)$. Thus
\eqref{eq:arith-giant-error}, \eqref{eq:arith-symbolic-error}, and
\eqref{eq:arith-bulk-error} remain negligible after the frequency and pattern
sums. These errors are uniform for each comparison of assignments. All
surviving Page multipliers are functions only of the real and $Q_f$
coordinates and are bounded by two per sampled prime.

\subsection{The signed comparison}

We prove \eqref{eq:src56}. Condition in the resulting main expression on
the real coordinates, the $Q_f$ coordinates, the nonbulk prime values, and
the top-giant residues. The bulk residues on $d$ are then independent
Haar units. The only remaining factor using them is
\eqref{eq:arith-spectator-product}. In particular the support descent,
\eqref{eq:src58}, and the integrated symbolic line factors have left no
additional bulk condition at the other small primes.

At a fixed $q\mid d$, each bulk slot is an independent Haar unit. In either
diagram its leaf products are marginally independent Haar units, because
the leaves use disjoint nonempty sets of slots. Jointly, their law is exactly
uniform on the subgroup specified by equality of the products in
corresponding overlap components. To verify this, regard each bulk slot as
an edge of the bipartite overlap multigraph. Its value contributes to the
product at both endpoints. In each connected component the product of the
left vertex products must equal the product of the right vertex products.
Conversely, prescribe vertex products satisfying this relation, set the
non-tree edges to one, and choose a spanning tree. Successively eliminate a
terminal vertex by assigning its incident tree edge the value needed for
that vertex. The last vertex is satisfied by the product relation. This
proves surjectivity onto exactly the asserted subgroup. A homomorphism of
finite groups sends the uniform distribution to the uniform distribution
on its image, proving the claimed joint law.

All frequencies and nonbulk constants needed for
Lemma~\ref{lem:tree-comparison} are units at $q$, and the ancestor and child
relations are those of \eqref{eq:src36}. The selected spectator primes
satisfy \eqref{eq:src35} uniformly, with a bound tending to zero; their
sizes tend to infinity uniformly as well. Therefore \eqref{eq:src38}
provides a bound $\epsilon_k(L)$, where $\epsilon_k(L)\to0$, for the
absolute correlation at every spectator, uniformly in the conditioned
data. By CRT the bulk coordinates at distinct spectators are independent,
so their combined bound is $\epsilon_k(L)^m$.

The remaining total costs, including the representative domination,
bounded masses, and all frequency sums, are $\exp(O_k(m))$. Consequently
\[
 \left|\E\sum_s\mathcal G K_{\rm ar}
                   A_l^{(1)}(s)\overline{A_l^{(2)}(s)}\right|
 \leq \epsilon_k(L)^m\exp(O_k(m))
       +O\bigl(\exp(-e^{.00125L})\bigr)
 =\exp(-\omega_k(m)),
\]
as required in \eqref{eq:src56}.

\subsection{The norm comparison and the frequency sums}

We finally prove \eqref{eq:src55}. Use the nonnegative upper bound described
above. At the $d$ coordinates, \eqref{eq:src37} and Cauchy--Schwarz bound the
expectation of the absolute product of two tree values by $3^r$ for each
spectator, hence by $3^{rm}$ in total. Retain the frequency divisibility
conditions, and discard the Page multipliers at cost $2^{rm+O(1)}$. The
bulk coordinates modulo $Q_f$ are now independent Haar units for the
upper bound.

Because the assignments agree, the bulk slots split into exactly the same
subtrees in the two histories. Fix the other necessary residues. Expose the
product of all bulk slots and then expose the products on successive subtree
splits, proceeding from the root downwards. Conditional on a parent product,
the product $M_+$ in one child is uniform on the unit group modulo $Q_f$,
and $M_-$ is determined by $M_+M_-$. This follows directly from the product
map on two disjoint collections of independent Haar units. Ancestor pivots
in both histories depend only on previously exposed splits and are known
with the precision supplied by \eqref{eq:arith-frequency-modulus}.

We evaluate support indicators under this original Haar measure, without
first conditioning on the validity of the complete histories. All frequencies
are fixed at the outset. A current giant's unit tests against descendant
frequencies therefore depend only on already exposed data. Tests on a newly
created pivot are evaluated after the split that creates it; tests on later
pivots are deferred to their own splits, or discarded for this nonnegative
upper bound. Thus the retained earlier tests do not select an unexposed split.

At a node, the support requires $X_\pm C_\pm$ to be units modulo its
frequency. If that unit test fails we stop that branch, whose contribution
is zero. Otherwise the condition
\begin{equation}\label{eq:arith-square-congruence}
 vX_-C_-M_-\equiv wX_+C_+M_+\pmod{|s|}
\end{equation}
requires $(v,s)=(w,s)$, since all the other factors are units. Put
$g=(v,w,s)$. In a soluble case, division by $g$ makes both frequency
coefficients units modulo
\begin{equation}\label{eq:arith-reduced-modulus}
                                a=\frac{|s|}{(v,w,s)}.
\end{equation}
Substituting the fixed parent product $M_+M_-$ into
\eqref{eq:arith-square-congruence} then fixes $M_+^2$ to a specified unit
modulo $a$. The other history at the same node fixes a square modulo
$a'=|s'|/(v',w',s')$.

For each odd prime power, a unit has at most two square roots, and for a
power of two it has at most four. Thus the simultaneous square equations
have at most $2^{\omega([a,a'])+1}$ solutions modulo $[a,a']$, where brackets
denote the least common multiple. If they are inconsistent, their
probability is zero. The ordinary divisor bound and
$n/\phi(n)\leq\tau(n)$ give, uniformly for the moduli at hand,
\[
 2^{\omega(n)+1}\frac{n}{\phi(n)}=\exp(o_k(m)),
             \qquad n\leq\exp(O_k(m)).
\]
It follows that the conditional probability at this split is at most
\begin{equation}\label{eq:arith-split-probability}
                          \frac{\exp(o_k(m))}{[a,a']}.
\end{equation}
Reduction of Haar units modulo $Q_f$ to this modulus is uniform. All
earlier tests are measurable with respect to the earlier exposed products,
so these conditional upper bounds multiply along the tree. Failed tests
contribute zero and do not alter the upper bound.

It remains to sum \eqref{eq:arith-split-probability} over internal
frequencies. The requisite numerical estimate is
\begin{equation}\label{eq:arith-lcm-sum}
                  \sum_{a,a'\leq V}\frac1{[a,a']}
                           \ll(1+\log V)^3.
\end{equation}
Indeed, write $a=du$, $a'=dv$ with $(u,v)=1$, and then discard the
coprimality condition. The left side is at most
\[
                  \sum_{d\leq V}\frac1d
                       \left(\sum_{u\leq V/d}\frac1u\right)^2,
\]
which proves \eqref{eq:arith-lcm-sum}.
With child frequencies $v,w$ fixed, every $s$ giving a specified
$a=|s|/(v,w,s)$ is of the form $s=\pm ad$ with $d\mid(v,w)$.
Its multiplicity is therefore at most $2\tau((v,w))$. The analogous bound
holds in the other history. Since these gcds and all the frequency bounds
are at most $\exp(O_k(m))$, \eqref{eq:arith-lcm-sum} shows that each pair of
internal frequencies sums to $\exp(o_k(m))$, uniformly in the fixed child
frequencies.

Drop the equality of the two root frequencies if necessary for this
nonnegative numerical bound. Sum the root pair first while holding the
children fixed, then proceed down the tree. Only the $2r$ leaf frequencies
remain, giving at most $(2V_0)^{2r}$ choices. Combining this with
\eqref{eq:arith-xi-bound} and $E_0=\Delta_0+\sqrt m$ yields
\begin{equation}\label{eq:arith-leaf-budget}
 (2V_0)^{2r}e^{-r\Delta_0+O_k(1)}
                         =\exp(r\Delta_0+o_k(m)).
\end{equation}
All the other losses fit $Crm+o_k(m)$ in the exponent. These are the
$3^{rm}$ spectator bound, bounded per-coordinate masses and Page factors,
the domination \eqref{eq:src59}, and $O_k(1)$ equality patterns. The
constant can be absolute: the internal-sample count is $O(kr)$, and
$kL/m\asymp k^{-3}$; fixed-$k$ constants not proportional to $m$ are
absorbed by $o_k(m)$. Inserting the negligible approximation errors gives
\[
                  \E\sum_s\mathcal G|A_l(s)|^2
                    \leq\exp\bigl(r(\Delta_0+Cm)+o_k(m)\bigr).
\]
This is \eqref{eq:src55}, and completes the proof of
Proposition~\ref{prop:comparisons}.
